\documentclass[10pt,aspectratio=169]{beamer}

% ------------------------------------------------------------------
% Tema 0: Eines per a una estadística computacional
% Grau en Enginyeria Informàtica - Estadística
% ------------------------------------------------------------------

\usepackage[utf8]{inputenc}
\usepackage[T1]{fontenc}
\usepackage[catalan]{babel}
\usepackage{amsmath,amssymb}
\usepackage{booktabs}
\usepackage{array}
\usepackage{tikz}
\usepackage{pgfplots}
\usepackage{xcolor}
\usepackage{listings}
\usepackage{hyperref}
\usetikzlibrary{positioning,arrows.meta,calc,shapes.geometric}
\pgfplotsset{compat=1.18}

% ------------------------------------------------------------------
% Estil visual coherent amb les presentacions anteriors
% ------------------------------------------------------------------
\definecolor{UABGreen}{RGB}{0,103,71}
\definecolor{UABLight}{RGB}{224,241,236}
\definecolor{DeepBlue}{RGB}{20,65,110}
\definecolor{AccentOrange}{RGB}{230,126,34}
\definecolor{SoftGray}{RGB}{245,247,250}
\definecolor{CodeBack}{RGB}{248,248,248}
\definecolor{CodeFrame}{RGB}{210,215,220}

\usetheme{Madrid}
\usecolortheme[named=UABGreen]{structure}
\setbeamercolor{title}{fg=white,bg=UABGreen}
\setbeamercolor{frametitle}{fg=white,bg=UABGreen}
\setbeamercolor{block title}{fg=white,bg=DeepBlue}
\setbeamercolor{block body}{fg=black,bg=SoftGray}
\setbeamercolor{alerted text}{fg=AccentOrange}
\setbeamertemplate{navigation symbols}{}
\setbeamertemplate{itemize item}{\color{UABGreen}$\blacktriangleright$}
\setbeamertemplate{itemize subitem}{\color{AccentOrange}$\bullet$}
\setbeamertemplate{enumerate item}{\color{UABGreen}\insertenumlabel.}

\setbeamertemplate{footline}{%
  \leavevmode%
  \hbox{%
  \begin{beamercolorbox}[wd=.72\paperwidth,ht=2.4ex,dp=1ex,leftskip=1.4em]{author in head/foot}%
    Estadística -- Grau en Enginyeria Informàtica
  \end{beamercolorbox}%
  \begin{beamercolorbox}[wd=.28\paperwidth,ht=2.4ex,dp=1ex,center]{date in head/foot}%
    \insertframenumber{} / \inserttotalframenumber
  \end{beamercolorbox}}%
  \vskip0pt%
}

\lstdefinelanguage{R}{%
  morekeywords={if,else,repeat,while,function,for,in,next,break,TRUE,FALSE,NULL,NA,NaN,Inf,
    library,install.packages,mean,median,sd,var,summary,table,prop.table,round,
    ggplot,aes,geom_point,geom_line,geom_col,geom_histogram,geom_boxplot,geom_smooth,
    labs,theme_minimal,summarise,group_by,mutate,read_csv,write_csv,filter,arrange,
    select,rename,set.seed,runif,rnorm,rbinom,rpois,sample,replicate,system.time,
    tibble,data.frame,head,str,glimpse,sessionInfo,dir.create,file.path,list.files,
    plan,multisession,future_map_dbl,map_dbl},
  sensitive=true,
  morecomment=[l]\#,
  morestring=[b]",
  morestring=[b]'
}
\lstdefinelanguage{bash}{%
  morekeywords={git,cd,mkdir,ls,pwd,quarto,Rscript,status,add,commit,push,pull,clone,init},
  sensitive=true,
  morecomment=[l]\#,
  morestring=[b]",
  morestring=[b]'
}
\lstset{%
  language=R,
  basicstyle=\ttfamily\scriptsize,
  keywordstyle=\color{DeepBlue}\bfseries,
  commentstyle=\color{gray},
  stringstyle=\color{AccentOrange},
  backgroundcolor=\color{CodeBack},
  frame=single,
  rulecolor=\color{CodeFrame},
  breaklines=true,
  columns=fullflexible,
  keepspaces=true,
  showstringspaces=false,
  literate={à}{{\`a}}1 {è}{{\`e}}1 {é}{{\'e}}1 {í}{{\'i}}1 {ï}{{\"i}}1 {ò}{{\`o}}1 {ó}{{\'o}}1 {ú}{{\'u}}1 {ü}{{\"u}}1 {ç}{{\c{c}}}1 {À}{{\`A}}1 {È}{{\`E}}1 {É}{{\'E}}1 {Í}{{\'I}}1 {Ò}{{\`O}}1 {Ó}{{\'O}}1 {Ú}{{\'U}}1
}

\newcommand{\R}{\textsf{R}}
\newcommand{\RStudio}{\textsf{RStudio/Posit}}
\newcommand{\GitHub}{\textsf{GitHub}}
\newcommand{\Quarto}{\textsf{Quarto}}
\newcommand{\Git}{\textsf{Git}}
\newcommand{\important}[1]{\textcolor{UABGreen}{\textbf{#1}}}
\newcommand{\exemple}[1]{\textcolor{AccentOrange}{\textbf{#1}}}

% ------------------------------------------------------------------
% Dades de la presentació
% ------------------------------------------------------------------
\title[Tema 0: eines computacionals]{Tema 0. Eines per a una estadística computacional}
\subtitle{R, RStudio/Posit, Quarto, GitHub i fluxos de treball reproduïbles}
\author{Estadística -- Grau en Enginyeria Informàtica}
\date{Curs 2026--27}

\begin{document}

% ------------------------------------------------------------------
\begin{frame}[plain]
  \titlepage
  \vspace{-0.4cm}
  \begin{center}
  \small
  \textcolor{DeepBlue}{De les dades al codi, del codi als resultats, dels resultats a una decisió}
  \end{center}
\end{frame}

% ------------------------------------------------------------------
\begin{frame}{Per què un tema 0?}
\begin{columns}[T]
\column{0.55\textwidth}
\begin{itemize}
  \item En aquesta assignatura no només farem càlculs: també aprendrem a \important{treballar amb dades}.
  \item L'objectiu és que cada resultat sigui:
  \begin{itemize}
    \item traçable,
    \item reproduïble,
    \item comunicable,
    \item i fàcil de revisar.
  \end{itemize}
  \item Això és especialment important en informàtica: dades, scripts, versions i informes han de conviure.
\end{itemize}

\column{0.40\textwidth}
\begin{block}{Idea central}
\centering
\large
\important{Estadística = dades + probabilitat + codi + interpretació}
\end{block}
\vspace{0.2cm}
\begin{alertblock}{Missatge pràctic}
El codi que no es pot reproduir és difícil de validar.
\end{alertblock}
\end{columns}
\end{frame}

% ------------------------------------------------------------------
\begin{frame}{Objectius de la sessió}
Al final d'aquest tema hauríeu de poder:
\begin{enumerate}
  \item Obrir un projecte d'\RStudio{} de manera ordenada.
  \item Crear i executar un script senzill en \R.
  \item Importar dades i generar una primera taula o gràfic.
  \item Preparar un informe breu amb \Quarto.
  \item Fer servir \Git{} i \GitHub{} per lliurar una pràctica reproduïble.
  \item Entendre per què la simulació i la computació paral·lela apareixeran al llarg del curs.
\end{enumerate}

\vspace{0.25cm}
\begin{block}{No és un curs de programació avançada}
És una introducció al \important{flux de treball computacional} que farem servir per aprendre estadística.
\end{block}
\end{frame}

% ------------------------------------------------------------------
\begin{frame}{Les eines del curs}
\begin{columns}[T]
\column{0.48\textwidth}
\begin{block}{Treball local}
\begin{itemize}
  \item \R: llenguatge per a càlcul estadístic.
  \item \RStudio: entorn de treball.
  \item \Quarto: informes amb text, codi i resultats.
\end{itemize}
\end{block}

\begin{block}{Treball amb versions}
\begin{itemize}
  \item \Git: historial de canvis.
  \item \GitHub: repositori remot i lliuraments.
\end{itemize}
\end{block}

\column{0.48\textwidth}
\begin{block}{Paquets d'\R}
\begin{itemize}
  \item \texttt{tidyverse}: manipulació i visualització.
  \item \texttt{readr}: lectura de fitxers.
  \item \texttt{ggplot2}: gràfics.
  \item \texttt{future}, \texttt{furrr}: paral·lelització.
\end{itemize}
\end{block}

\begin{alertblock}{Durant el curs}
No cal memoritzar totes les funcions: cal aprendre a construir un procés net i verificable.
\end{alertblock}
\end{columns}
\end{frame}

% ------------------------------------------------------------------
\begin{frame}{Flux de treball recomanat}
\begin{center}
\begin{tikzpicture}[node distance=1.15cm, every node/.style={font=\small}]
\node[draw, rounded corners, fill=UABLight, minimum width=2.3cm, minimum height=0.8cm] (data) {Dades};
\node[draw, rounded corners, fill=SoftGray, right=of data, minimum width=2.3cm, minimum height=0.8cm] (code) {Codi};
\node[draw, rounded corners, fill=SoftGray, right=of code, minimum width=2.3cm, minimum height=0.8cm] (results) {Resultats};
\node[draw, rounded corners, fill=SoftGray, right=of results, minimum width=2.3cm, minimum height=0.8cm] (report) {Informe};
\node[draw, rounded corners, fill=UABLight, below=of code, minimum width=2.6cm, minimum height=0.8cm] (git) {Git/GitHub};
\draw[-{Latex[length=3mm]}, thick, color=UABGreen] (data) -- (code);
\draw[-{Latex[length=3mm]}, thick, color=UABGreen] (code) -- (results);
\draw[-{Latex[length=3mm]}, thick, color=UABGreen] (results) -- (report);
\draw[-{Latex[length=3mm]}, thick, color=AccentOrange] (git) -- (code);
\draw[-{Latex[length=3mm]}, thick, color=AccentOrange] (report) -- ++(0,-1) -| (git);
\end{tikzpicture}
\end{center}

\vspace{0.25cm}
\begin{block}{Principi de reproduïbilitat}
Una altra persona hauria de poder obtenir els mateixos resultats executant el vostre codi en el mateix projecte.
\end{block}
\end{frame}

% ------------------------------------------------------------------
\begin{frame}[fragile]{R com a calculadora estadística}
\begin{columns}[T]
\column{0.54\textwidth}
\begin{lstlisting}
# Operacions bàsiques
2 + 3
sqrt(16)
log(100)

# Objectes
x <- c(12, 15, 13, 18, 20)
mean(x)
sd(x)
summary(x)
\end{lstlisting}

\column{0.42\textwidth}
\begin{itemize}
  \item En \R{} assignem valors amb \texttt{<-}.
  \item Un vector és una col·lecció de valors.
  \item Moltes funcions estadístiques bàsiques ja estan incorporades.
\end{itemize}

\begin{alertblock}{Bon hàbit}
Comenteu el codi amb \texttt{\#}: ajuda a entendre què esteu fent.
\end{alertblock}
\end{columns}
\end{frame}

% ------------------------------------------------------------------
\begin{frame}[fragile]{Primer exemple: temps d'execució}
\begin{columns}[T]
\column{0.57\textwidth}
\begin{lstlisting}
set.seed(123)

temps_ms <- c(31, 29, 34, 30, 35, 33, 32, 80)

mean(temps_ms)
median(temps_ms)
sd(temps_ms)
summary(temps_ms)
\end{lstlisting}

\column{0.39\textwidth}
\begin{block}{Pregunta estadística}
Quin és un temps típic d'execució? Hi ha algun valor sospitós?
\end{block}

\begin{itemize}
  \item La mitjana és sensible a valors extrems.
  \item La mediana pot ser més robusta.
  \item La desviació típica mesura variabilitat.
\end{itemize}
\end{columns}
\end{frame}

% ------------------------------------------------------------------
\begin{frame}[fragile]{Dades en format taula}
\begin{columns}[T]
\column{0.55\textwidth}
\begin{lstlisting}
library(tidyverse)

bench <- tibble(
  algorisme = c("A", "A", "A", "B", "B", "B"),
  mida = c(1000, 5000, 10000, 1000, 5000, 10000),
  temps_ms = c(8, 40, 85, 6, 29, 62)
)

bench
\end{lstlisting}

\column{0.41\textwidth}
\begin{block}{Una fila = una observació}
Cada fila representa una execució o mesura.
\end{block}

\begin{block}{Una columna = una variable}
Per exemple: algorisme, mida de l'entrada, temps d'execució.
\end{block}
\end{columns}
\end{frame}

% ------------------------------------------------------------------
\begin{frame}[fragile]{Una primera visualització}
\begin{columns}[T]
\column{0.58\textwidth}
\begin{lstlisting}
ggplot(bench, aes(x = mida, y = temps_ms,
                  color = algorisme)) +
  geom_point(size = 3) +
  geom_line() +
  labs(
    x = "Mida de l'entrada",
    y = "Temps d'execució (ms)",
    color = "Algorisme"
  ) +
  theme_minimal()
\end{lstlisting}

\column{0.38\textwidth}
\begin{itemize}
  \item Els gràfics ajuden a detectar patrons.
  \item En informàtica són útils per estudiar rendiment.
  \item El codi genera el gràfic: no cal copiar-lo manualment.
\end{itemize}

\begin{alertblock}{Connexió amb el curs}
La descriptiva bivariant començarà precisament amb relacions entre dues variables.
\end{alertblock}
\end{columns}
\end{frame}

% ------------------------------------------------------------------
\begin{frame}{Projectes d'RStudio}
\begin{columns}[T]
\column{0.52\textwidth}
\begin{block}{Què és un projecte?}
Una carpeta que conté tot el necessari per a una anàlisi:
\begin{itemize}
  \item dades,
  \item scripts,
  \item informes,
  \item resultats,
  \item historial de versions.
\end{itemize}
\end{block}

\column{0.44\textwidth}
\begin{block}{Estructura recomanada}
\ttfamily\small
estadistica-ei/\\
\quad data/\\
\quad R/\\
\quad reports/\\
\quad figures/\\
\quad README.md\\
\quad practica0.qmd
\end{block}
\end{columns}

\vspace{0.25cm}
\begin{alertblock}{Regla pràctica}
Eviteu rutes absolutes com \texttt{/home/usuari/Escriptori/...}. Feu servir camins relatius dins del projecte.
\end{alertblock}
\end{frame}

% ------------------------------------------------------------------
\begin{frame}[fragile]{Llegir dades}
\begin{columns}[T]
\column{0.55\textwidth}
\begin{lstlisting}
library(tidyverse)

logs <- read_csv("data/logs_servidor.csv")

glimpse(logs)
summary(logs)
\end{lstlisting}

\column{0.41\textwidth}
\begin{block}{Exemples de dades del curs}
\begin{itemize}
  \item Logs de servidor.
  \item Temps d'execució.
  \item Resultats de simulacions.
  \item Errors de transmissió.
  \item Peticions per segon.
\end{itemize}
\end{block}

\begin{alertblock}{Important}
No modifiqueu manualment les dades originals: creeu scripts que facin la neteja.
\end{alertblock}
\end{columns}
\end{frame}

% ------------------------------------------------------------------
\begin{frame}{Quarto: informes reproduïbles}
\begin{columns}[T]
\column{0.50\textwidth}
\begin{block}{Un informe Quarto combina}
\begin{itemize}
  \item text explicatiu,
  \item codi en \R,
  \item taules,
  \item gràfics,
  \item conclusions.
\end{itemize}
\end{block}

\column{0.46\textwidth}
\begin{alertblock}{Per què és útil?}
Si canvien les dades o el codi, l'informe es pot tornar a generar automàticament.
\end{alertblock}

\begin{block}{Lliuraments}
Les pràctiques poden contenir el fitxer \texttt{.qmd}, el codi, les dades i el resultat en HTML o PDF.
\end{block}
\end{columns}
\end{frame}

% ------------------------------------------------------------------
\begin{frame}[fragile]{Un fragment mínim de Quarto}
\begin{lstlisting}[language=bash]
---
title: "Pràctica 0"
author: "Nom Cognoms"
format: html
---

## Objectiu
Analitzem els temps d'execució d'un algoritme.

```{r}
temps <- c(31, 29, 34, 30, 35, 33, 32, 80)
summary(temps)
```

## Comentari
La mediana és menys sensible al valor extrem que la mitjana.
\end{lstlisting}
\end{frame}

% ------------------------------------------------------------------
\begin{frame}{Git: control de versions}
\begin{columns}[T]
\column{0.52\textwidth}
\begin{block}{Què resol Git?}
\begin{itemize}
  \item Guardar l'historial del projecte.
  \item Recuperar versions anteriors.
  \item Saber què ha canviat.
  \item Treballar de forma ordenada.
\end{itemize}
\end{block}

\column{0.44\textwidth}
\begin{alertblock}{Sense Git}
\ttfamily\small
practica.R\\
practica\_final.R\\
practica\_final\_bona.R\\
practica\_final\_bona2.R
\end{alertblock}

\begin{block}{Amb Git}
Un fitxer, amb historial de canvis.
\end{block}
\end{columns}
\end{frame}

% ------------------------------------------------------------------
\begin{frame}[fragile]{Cicle bàsic amb Git}
\begin{columns}[T]
\column{0.54\textwidth}
\begin{lstlisting}[language=bash]
git status

git add practica0.qmd R/analisi.R

git commit -m "Afegeix primera anàlisi descriptiva"

git push
\end{lstlisting}

\column{0.42\textwidth}
\begin{enumerate}
  \item Miro l'estat del projecte.
  \item Marco els fitxers que vull guardar.
  \item Faig una fotografia del projecte.
  \item La pujo a GitHub.
\end{enumerate}

\begin{alertblock}{Missatges de commit}
Han de descriure el canvi, no només dir ``canvis''.
\end{alertblock}
\end{columns}
\end{frame}

% ------------------------------------------------------------------
\begin{frame}{GitHub: lliuraments i col·laboració}
\begin{columns}[T]
\column{0.50\textwidth}
\begin{block}{Què aportarà GitHub al curs?}
\begin{itemize}
  \item Repositoris per a pràctiques.
  \item Historial de lliuraments.
  \item README amb instruccions.
  \item Possibilitat de revisar codi.
\end{itemize}
\end{block}

\column{0.46\textwidth}
\begin{block}{Un bon repositori conté}
\begin{itemize}
  \item \texttt{README.md}: què fa el projecte.
  \item \texttt{data/}: dades necessàries.
  \item \texttt{R/}: scripts.
  \item \texttt{reports/}: informe final.
  \item commits clars i ordenats.
\end{itemize}
\end{block}
\end{columns}
\end{frame}

% ------------------------------------------------------------------
\begin{frame}[fragile]{Exemple de README}
\begin{lstlisting}[language=bash]
# Pràctica 0: temps d'execució

Objectiu: analitzar descriptivament els temps d'execució
simulats de dos algoritmes.

## Fitxers
- data/temps_algoritmes.csv: dades d'entrada
- R/analisi.R: codi principal
- reports/practica0.qmd: informe reproduïble

## Com reproduir els resultats
1. Obrir el projecte a RStudio.
2. Executar reports/practica0.qmd.
\end{lstlisting}
\end{frame}

% ------------------------------------------------------------------
\begin{frame}[fragile]{Simulació: una idea que apareixerà sovint}
\begin{columns}[T]
\column{0.56\textwidth}
\begin{lstlisting}
set.seed(2026)

# Simulem 1000 temps de resposta en ms
temps <- rnorm(n = 1000, mean = 120, sd = 15)

mean(temps)
quantile(temps, c(.05, .50, .95))
\end{lstlisting}

\column{0.40\textwidth}
\begin{itemize}
  \item La simulació ens ajuda a entendre fenòmens aleatoris.
  \item Permet experimentar sense tenir encara dades reals.
  \item Serà central en probabilitat, inferència i bootstrap.
\end{itemize}

\begin{alertblock}{Reproduïbilitat}
\texttt{set.seed()} fa que una simulació es pugui repetir.
\end{alertblock}
\end{columns}
\end{frame}

% ------------------------------------------------------------------
\begin{frame}[fragile]{Exemple: simular peticions a un servidor}
\begin{columns}[T]
\column{0.58\textwidth}
\begin{lstlisting}
set.seed(123)

# Nombre de peticions per segon
peticions <- rpois(n = 60, lambda = 2.5)

tibble(segon = 1:60, peticions = peticions) |>
  ggplot(aes(x = segon, y = peticions)) +
  geom_col() +
  theme_minimal()
\end{lstlisting}

\column{0.38\textwidth}
\begin{block}{Pregunta}
Quina és la probabilitat de tenir més de 5 peticions en un segon?
\end{block}

\begin{itemize}
  \item Aquest exemple connecta amb la distribució de Poisson.
  \item També connecta amb capacitat, rendiment i cues.
\end{itemize}
\end{columns}
\end{frame}

% ------------------------------------------------------------------
\begin{frame}{Per què parlar de computació paral·lela?}
\begin{columns}[T]
\column{0.52\textwidth}
\begin{block}{Moltes tasques estadístiques són repetitives}
\begin{itemize}
  \item 10.000 simulacions Monte Carlo.
  \item 5.000 mostres bootstrap.
  \item Comparar molts paràmetres.
  \item Avaluar molts escenaris.
\end{itemize}
\end{block}

\column{0.44\textwidth}
\begin{alertblock}{Idea}
Si les repeticions són independents, sovint podem repartir-les entre diversos nuclis del processador.
\end{alertblock}

\begin{block}{En aquest curs}
Ho veurem només de manera introductòria i aplicada.
\end{block}
\end{columns}
\end{frame}

% ------------------------------------------------------------------
\begin{frame}[fragile]{Primer tast de paral·lelització amb R}
\begin{lstlisting}
library(future)
library(furrr)

plan(multisession, workers = 4)

simula_mitjana <- function(i) {
  x <- rnorm(10000, mean = 0, sd = 1)
  mean(x)
}

resultats <- future_map_dbl(1:1000, simula_mitjana)

mean(resultats)
sd(resultats)
\end{lstlisting}

\begin{block}{Lectura del codi}
Fem 1000 simulacions independents i les distribuïm entre diversos processos.
\end{block}
\end{frame}

% ------------------------------------------------------------------
\begin{frame}{Bones pràctiques des del primer dia}
\begin{columns}[T]
\column{0.50\textwidth}
\begin{block}{Codi}
\begin{itemize}
  \item Noms de variables clars.
  \item Comentaris breus i útils.
  \item Scripts ordenats per blocs.
  \item Cap resultat copiat manualment.
\end{itemize}
\end{block}

\column{0.46\textwidth}
\begin{block}{Projecte}
\begin{itemize}
  \item Carpeta estructurada.
  \item Fitxers amb noms informatius.
  \item README actualitzat.
  \item Commits freqüents i descriptius.
\end{itemize}
\end{block}
\end{columns}

\vspace{0.25cm}
\begin{alertblock}{Objectiu}
Que una altra persona pugui entendre què heu fet sense preguntar-vos-ho.
\end{alertblock}
\end{frame}

% ------------------------------------------------------------------
\begin{frame}{Errors habituals}
\begin{columns}[T]
\column{0.48\textwidth}
\begin{alertblock}{Eviteu}
\begin{itemize}
  \item Codi només a la consola.
  \item Fitxers sense estructura.
  \item Rutes absolutes.
  \item Gràfics enganxats manualment.
  \item No fer commits fins al final.
\end{itemize}
\end{alertblock}

\column{0.48\textwidth}
\begin{block}{Feu}
\begin{itemize}
  \item Scripts i informes reproduïbles.
  \item Projectes d'RStudio.
  \item Camins relatius.
  \item Taules i gràfics generats amb codi.
  \item Commits petits i freqüents.
\end{itemize}
\end{block}
\end{columns}
\end{frame}

% ------------------------------------------------------------------
\begin{frame}{Per practicar}
\begin{block}{Tasques}
\begin{enumerate}
  \item Instal·lar o comprovar \R, \RStudio{}, \Git{} i \Quarto.
  \item Crear un repositori a \GitHub{} per a l'assignatura.
  \item Crear un projecte d'\RStudio{} connectat al repositori.
  \item Escriure un script que simuli o llegeixi dades de temps d'execució.
  \item Crear un informe \texttt{practica0.qmd} amb una taula, un gràfic i una interpretació breu.
  \item Fer alguns commits i pujar-los a \GitHub{}.
\end{enumerate}
\end{block}
\end{frame}


\end{document}
